\documentclass{article}

\usepackage[utf8]{inputenc}  % Para caracteres especiais
\usepackage[brazil]{babel} % Para a língua portuguesa

\usepackage[a4paper,top=2cm,bottom=2cm,left=2cm,right=2cm,marginparwidth=1.75cm]{geometry}

% Pacotes úteis
\usepackage{amsmath}
\usepackage{graphicx}
\usepackage[colorlinks=true, allcolors=blue]{hyperref}
\usepackage[alf]{abntex2cite} % Referências no padrão ABNT, conforme o edital

\title{Personalização de modelos com poucos dados para detecção de risco de hipoglicemia em diabetes tipo 1}
\author{[Nome completo do(a) candidato(a)]}
\date{\today}

\begin{document}
\maketitle

% O Abstract não é necessário
% \begin{abstract}
% Your abstract.
% \end{abstract}

\section*{Uso de IA Generativa}

% TODO: manter apenas a declaração verdadeira para a versão submetida.
% O edital exige que qualquer uso de IA generativa seja declarado.
% Opção 1: Declaro que nenhuma ferramenta de inteligência artificial generativa foi utilizada na elaboração desta proposta.
% Opção 2: Declaro que a ferramenta [nome] foi utilizada para [partes da proposta e finalidade].
\textbf{[TODO: inserir a declaração de uso de IA generativa aplicável.]}

\section{Introdução}

% TODO: apresentar o problema e a pergunta de pesquisa, contextualizar sua
% relevância e indicar concretamente como será investigado. O objetivo da
% pesquisa não deve ser descrito como apenas desenvolver um produto.
\textbf{[TODO: redigir a introdução.]}

\section{Fundamentação teórica / trabalhos relacionados}

% TODO: sintetizar os conceitos específicos e os trabalhos relacionados,
% sustentando as afirmações com referências atuais e relevantes.
\textbf{[TODO: redigir a fundamentação teórica e os trabalhos relacionados.]}

\section{Metodologia}

% TODO: descrever o desenho da pesquisa, dados, procedimentos, métodos,
% comparação e métricas de avaliação. Incluir figuras se forem úteis.
\textbf{[TODO: redigir a metodologia.]}

\section{Resultados esperados}

% TODO: explicitar resultados verificáveis e contribuições esperadas,
% evitando promessas que não possam ser avaliadas.
\textbf{[TODO: descrever os resultados esperados.]}

\bibliographystyle{abntex2-alf}
\bibliography{referencias}

\end{document}
